\section{Related Work}

\textbf{Benchmarks for autonomous research agents.}
Language-model agents are increasingly evaluated on executable research and engineering tasks. 
Early benchmarks such as MLAgentBench~\citep{MLAgentBench2024Huang}, MLE-bench~\citep{chan2025mlebench}, and RE-Bench~\citep{REbench2025Wijk} require agents to modify code, run experiments, and improve machine-learning systems under realistic constraints. 
More recent benchmarks, including PostTrainBench~\citep{posttrainbench_2026}, MLS-Bench~\citep{lyu2026mlsbench}, and AutoLab~\citep{xu2026autolab}, extend this paradigm to longer-horizon, resource-bounded settings in which agents repeatedly propose modifications, observe empirical feedback, and refine executable artifacts. 
Frontier-Eng~\citep{chi2026frontiereng} and FML-bench~\citep{zou2026fmlbench} further analyze aggregate search behaviors such as improvement frequency, exploration diversity, and search depth. However, these benchmarks primarily evaluate final performance or global properties of the search trajectory. 
Our work instead jointly evaluates process competence and experience-driven self-improvement, moving beyond outcome scores to reveal both why agents succeed or fail and whether they improve through experience.

\textbf{Process-level evaluation of long-horizon agents.}
Several studies have moved beyond terminal success metrics to evaluate intermediate agent behavior. 
AgentBoard~\citep{AgentBoard2024Ma} introduces progress rate to measure advancement toward intermediate subgoals, while TRAJECT-Bench~\citep{he2026trajectbench} evaluates the correctness of tool selection, arguments, and execution order. 
WebStep~\citep{chung2026webstep} tracks semantic environment states to separate exploration reach from execution accuracy, whereas AgentLens~\citep{sahoo2026agentlens} compares software-engineering trajectories against successful process references to identify inefficient behaviors and ``lucky passes.'' 
These works analyze task execution through predefined subgoals, expected action structures, or instrumented process representations. 
Our work extends process-level evaluation to auto research workflows, providing objective, judge-free attribution without assuming canonical solution paths.

\textbf{Experience reuse and self-improving agents.}
Prior work has explored how agents can improve by retaining and reusing past experience. 
Reflexion~\citep{Shinn2023Reflexion} converts task feedback into verbal reflections that guide subsequent attempts, while ExpeL~\citep{Zhao2024ExpeL} extracts reusable insights from prior trajectories for cross-task transfer. 
LifelongAgentBench~\citep{zheng2025lifelongagentbench} and SEA-Eval~\citep{jiang2026seaeval} extend evaluation from isolated episodes to sequential task streams, measuring experience accumulation, skill transfer, and longer-term evolution. 
More recent evaluations, including SkillsBench~\citep{li2026skillsbench} and EvoAgentBench~\citep{gao2026evoagentbench}, study whether procedural knowledge or trace-derived abilities improve performance across tasks, showing that experience reuse can be beneficial but is often unstable and may cause negative transfer. 
Using an alternative evaluation design, concurrent work on EdgeBench~\citep{zhu2026edgebenchunveilingscalinglaws} studies scaling laws for learning from environments, a concept closely related to our intra-task self-improvement.
Our work further evaluates experience-driven self-improvement both within and across tasks in long-horizon auto research.

\textbf{Impact of Agent Harnesses.}
Agent performance depends not only on the underlying model but also on the harness that governs tool use, context construction, execution, and feedback. 
SWE-agent~\citep{Yang2024SWEagent} demonstrates that agent--computer interface design can substantially affect software-engineering performance. 
The Holistic Agent Leaderboard~\citep{kapoor2026holistic} jointly analyzes models, scaffolds, and benchmarks under standardized evaluation, while Harness-Bench~\citep{yao2026harnessbench} directly measures harness effects across multiple model backends and execution configurations. 
\citet{zhang2026stopcomparingllmagents} further argue that long-horizon agent comparisons require explicit harness disclosure and controlled evaluation protocols. 
Following this line, we fix the harness in our main cross-model comparison and use a native-harness ablation to assess the sensitivity of the results.
